\begin{abstract}
This dissertation investigates how evaluation techniques for retrieval-augmented generation (RAG) systems can be organized and applied to the analysis of document retrieval, answer generation, and source verifiability in a system for querying speeches from the 56th Legislature of the Brazilian Federal Senate. The applied, exploratory, and descriptive study adopts a multimethod approach and is informed by Design Science Research. The corpus comprises 15,729 speech records, of which 15,726 contain a full text or a usable summary. For the feasibility stage, a battery of 200 questions across 16 categories was developed, and a stratified pilot of 32 items was selected. In the reported run, all items were processed, and the automated judge assigned a mean score of 9.094 on a 0-to-10 scale using an exploratory rubric. The run also flagged a high risk of author mixing in 28 items. This score is not an independent human assessment: \texttt{gpt-5-nano} was used both to generate and to judge the answers. The main experiment has not yet been conducted; there is no validated reference set, controlled comparison among retrieval methods, or independent human evaluation. The available results therefore indicate operational feasibility of the workflow and inform protocol revision, but do not support conclusions about retrieval or generation performance or about the study hypotheses.
\end{abstract}
